\documentclass{situationist-preprint}

\title{The Recurrence of the Crisis}
\msid{PHE-17-a}
\msfamily{PHE-17}
\msdiscipline{Phenomenology}
\mstype{Principal manuscript}
\mskeywords{life-world; Husserl; Crisis; Galileo; mathematization; sedimentation; reactivation; statistics; computational science; scientific realism}

\begin{document}
\maketitle

\begin{abstract}
In the \emph{Crisis of the European Sciences}, Husserl argued that modern science is founded on the life-world, the world of pre-scientific experience, and that since Galileo it has forgotten this foundation, substituting its mathematical idealizations for the world they were devised to describe. Is the life-world prior to science, and was the crisis a single historical event? We argue that the life-world is prior to science in the order of sense and evidence, that science is a founded extension of it, and that such extensions tend to forget their foundation because their methods are transmitted as sedimented techniques that can be applied without being reactivated. The crisis Husserl diagnosed is therefore recurrent rather than unique: it arises wherever a successful method is technized and transmitted across generations of practitioners. We support the recurrence thesis with three episodes, Galilean mathematization, the institutionalization of statistical inference and the rise of computational science, and show that the account is compatible with scientific realism. The priority of the life-world concerns where the sense of scientific concepts comes from, not whether scientific theories are true; the forgetting is a misunderstanding of what the theories are, not a defect in what they say.
\end{abstract}

\section{The problem}

\subsection{Husserl's diagnosis}

Husserl's last work opens with a paradox. The European sciences are, by every internal standard, flourishing: their methods are rigorous, their results secure and their technical applications unprecedented. Yet they are in crisis, because they have lost their meaning for human life \citep[§§1--7]{Husserl1954}. The positive sciences exclude in principle the questions that matter most to human beings, about reason and unreason, freedom and meaning, and they have come to regard the world as nothing but the object of their methods. The diagnosis is not that science has made errors but that it has misunderstood itself.

The core of the diagnosis is the long §9, on Galileo's mathematization of nature. Galileo inherited geometry as a finished tradition, a technique of idealization that transformed the shapes of the surrounding world, which are vague and approximate, into the exact limit-shapes of geometrical space. He extended the technique from shapes to the qualities of bodies, the colors, sounds and warmths that have no geometrical form, by treating them as indices of underlying quantitative processes. Nature as a whole became a mathematical manifold, and the task of science became the construction of formulae that predict the course of experience. Husserl calls Galileo \enquote{at once a discovering and a concealing genius} \citep[§9h]{Husserl1954}: discovering, because the method was immensely fruitful; concealing, because it hid the origin of its own sense. The mathematized nature was taken for true nature, and the world of experience, from which the idealizations were drawn and in which their predictions are tested, was reduced to mere subjective appearance. Husserl calls this the \enquote{surreptitious substitution} of the mathematically substructed world of idealities for the only real world, the world given in perception, the everyday life-world \citep[§9h]{Husserl1954}.

Two theses must be distinguished in this diagnosis. One is the \emph{priority thesis}: the life-world is prior to science, in the sense that science is a founded extension of it, whose concepts derive their sense from life-world practices of measuring, counting and shaping, and whose results are confirmed, finally, by observations made in it. The other is the \emph{forgetting thesis}: science tends to forget this foundation and to take its own constructions for the whole of reality. A third question, which Husserl's text invites but does not settle, concerns history. Husserl tells the story of the forgetting as the story of modern European science, beginning with Galileo and culminating in the positivism of his own time. Is the crisis a unique event in the history of reason, or a recurrent structure?

\subsection{Reception}

The early phenomenological reception focused on the priority thesis. \citet{Gurwitsch1974} reconstructed the Galileo section as a theory of the constitution of physical science, showing in detail how the concepts of mathematical physics arise by idealization from perceptual concepts, and argued that the priority of the life-world is compatible with the objectivity of physics, since what is founded on the life-world is the sense of physical concepts, not their truth. \citet{Heelan1983} extended the analysis to the hermeneutics of scientific observation, arguing that instruments extend perception rather than replace it.

\citet{Carr1974} focused on the historical dimension. He argued that Husserl's turn to history in the \emph{Crisis} is not a departure from transcendental philosophy but a consequence of it: the sense of scientific concepts is constituted through a history of sedimentation, and transcendental reflection must become historical reflection to recover it. Carr also noted the tension between this view and Husserl's teleological philosophy of history, on which European history is the unfolding of an idea of reason that the crisis threatens. If the crisis is the crisis of that teleology, it is unique; if it is a consequence of the structure of sedimentation, it should recur.

\citet{Blumenberg1981} gave the most influential treatment of the forgetting thesis. In his essay on the life-world and technization, he argued that Husserl's notion of technization, the transformation of a method into a technique that can be practised without insight into its sense, describes a structural feature of all rational practice rather than a historical aberration. Techniques exist precisely so that their users need not reactivate their foundations; the forgetting is the condition of their efficiency. Blumenberg's later lectures on the theory of the life-world developed the point into a general account of the life-world as the domain of the self-evident, from which theory departs and to which it cannot return \citep{Blumenberg2010}.

More recently, \citet{Moran2012} has given a comprehensive commentary on the \emph{Crisis} that places the Galileo section within Husserl's account of the history of transcendental philosophy, and that stresses the ambiguity of the life-world concept: it is sometimes the concrete, culturally specific world of a community, and sometimes the invariant structure that all such worlds share. In analytic philosophy, \citet{Sellars1962} drew a related distinction between the manifest and the scientific images of man in the world, and argued that the scientific image is methodologically dependent on the manifest image but ontologically primary. Sellars's position is close to the one we shall defend on realism, though our account of the priority of the life-world is phenomenological, not Sellarsian.

\subsection{The thesis}

We argue for three claims. First, the life-world is prior to science: science is a founded extension of the life-world, and the priority is one of sense and evidence. Second, founded extensions tend to forget their foundation, because the sense of their concepts is transmitted in sedimented form and need not be reactivated by those who use them. Third, the crisis Husserl diagnosed is recurrent rather than unique: it arises wherever the conditions of forgetting obtain, and these conditions are generic features of successful sciences. We support the third claim by three episodes, and we show that the account is compatible with scientific realism.

\section{Framework}\label{s:framework}

\begin{definition}\label{d:lifeworld}
The \emph{life-world} is the world as given in pre-scientific experience: the world of perceived things, practical activities and intersubjective communication, within which every activity, including science, takes place.
\end{definition}

The life-world, so defined, is not a part of the world but the world itself as the horizon of all experience. PHE-11 argues that the world is the horizon of all horizons and has no horizon of its own; the life-world is that horizon as it is lived pre-theoretically, and it is never given as an object, though its structures can be described.

\begin{definition}\label{d:founded}
A practice $S$ is a \emph{founded extension} of the life-world if (i) the basic concepts of $S$ have their sense through operations of idealization, quantification or formalization performed on life-world practices, and (ii) the evidence for the results of $S$ is finally provided by experiences in the life-world: observations, readings of instruments, the perception of written records.
\end{definition}

\begin{definition}\label{d:sedimentation}
A concept or method is \emph{sedimented} when it is transmitted in a form that allows it to be used correctly without the user's re-performing the operations through which its sense was originally constituted. It is \emph{reactivated} when a user re-performs those operations and so makes its original sense evident again.
\end{definition}

The terms are Husserl's, from the \emph{Origin of Geometry} \citep{Husserl1939b}. Geometry, Husserl argues, is constituted by the first geometers through idealizing operations performed on the practices of surveying and shaping. Its results are fixed in writing, which allows them to be transmitted beyond the lifetime of their discoverers. The written proposition can be understood passively, as a meaningful sentence, without its sense being reactivated, and new propositions can be derived from it by the logical manipulation of written formulae. A tradition can thus grow by building on sedimented results whose original sense no one any longer reactivates. Husserl insists that the capacity for reactivation is in principle preserved, but that in practice it is lost, and that the loss is the condition of the tradition's growth \citep[see also][]{Derrida1962}.

\begin{definition}\label{d:forgetting}
A founded extension \emph{forgets its foundation} when its practitioners take its idealized objects for the whole of reality, and the life-world from which their sense derives for a mere appearance of them.
\end{definition}

\begin{definition}\label{d:crisis}
A \emph{crisis} of a founded extension is a state in which it has forgotten its foundation and, as a result, can no longer account for the relevance of its results to the life-world from which its questions arise.
\end{definition}

Definition~\ref{d:crisis} is narrower than Husserl's own notion, which includes the loss of meaning of science for human life as a whole. It retains the structural core: the crisis is not a failure of method or a crisis of results but a loss of self-understanding that becomes visible when the science cannot say what its results mean for the questions that occasioned them.

\begin{principle}\label{p:idealization}
Every science has basic concepts constituted by idealization, quantification or formalization of life-world practices, and confirms its results through observations made in the life-world.
\end{principle}

\begin{principle}\label{p:technization}
Every successful method tends to be technized: transmitted as a technique whose correct use does not require the reactivation of its sense.
\end{principle}

\begin{principle}\label{p:transmission}
A technized method that is transmitted across generations of practitioners is, after some generations, used by practitioners most of whom cannot reactivate its sense.
\end{principle}

Principle~\ref{p:idealization} is the content of the priority thesis, stated for all sciences. Principle~\ref{p:technization} is Husserl's thesis of technization \citep[§9g]{Husserl1954}, generalized by \citet{Blumenberg1981}: efficiency requires that the user of a technique not be required to understand its foundation. Principle~\ref{p:transmission} is a claim about the economy of training. Training in a science must be finite, while the history of its methods grows; at some point the methods are taught as techniques, not reconstructed from their origins.

\section{The argument}\label{s:argument}

\begin{theorem}\label{t:priority}
The life-world is prior to science: every science is a founded extension of the life-world.
\end{theorem}

\begin{proof}
By Principle~\ref{p:idealization}, every science has basic concepts constituted by idealization, quantification or formalization of life-world practices, which is condition (i) of Definition~\ref{d:founded}, and confirms its results through observations made in the life-world, which is condition (ii).
\end{proof}

The priority established by Theorem~\ref{t:priority} is a priority of sense and evidence. It is not a priority of being: the theorem does not say that the life-world is more real than the entities posited by science, nor that scientific entities are constructions out of life-world objects. That a concept has its sense through idealization does not prevent it from applying to something that exists independently of the idealization, just as the sense of the concept of a meridian derives from surveying practices while the earth's shape does not. We return to this point in Section~\ref{s:realism}.

\begin{proposition}\label{p:tendency}
Every successful science tends to forget its foundation.
\end{proposition}

\begin{proof}
By Principles~\ref{p:technization} and~\ref{p:transmission}, the methods of a successful science come to be used by practitioners most of whom cannot reactivate their sense. Such practitioners encounter the idealized objects of the science not as products of idealization but as given: the mathematical manifold, the probability distribution, the model, are for them simply what the science is about. Their access to the life-world from which the sense of these objects derives is the access of everyone else, pre-scientific and non-thematic. Nothing in their practice holds the two together, and the most economical self-interpretation available to them is that the objects of the science are the reality and the life-world a confused appearance of it. Where this self-interpretation prevails, the science has forgotten its foundation in the sense of Definition~\ref{d:forgetting}.
\end{proof}

Proposition~\ref{p:tendency} states a tendency, not a law. Some practitioners reactivate; some sciences institutionalize reflection on their foundations; and the self-interpretation described may be resisted by a culture that keeps the life-world in view. But the tendency is structural: it follows from the conditions under which a science grows, not from the failings of particular scientists.

\begin{theorem}\label{t:recurrence}
Crises in the sense of Definition~\ref{d:crisis} are recurrent: they tend to arise in every successful founded extension whose methods are technized and transmitted, and not only in the mathematical natural science of the modern period.
\end{theorem}

\begin{proof}
By Theorem~\ref{t:priority}, every science is a founded extension. By Proposition~\ref{p:tendency}, every successful science tends to forget its foundation. A science that has forgotten its foundation cannot account for the relevance of its results to the life-world, since such an account requires tracing the sense of its concepts back to the life-world practices from which they derive, which is what forgetting precludes. So every successful science tends towards crisis. Since there are many successful sciences, and since the conditions of the proof are not specific to any period, crises tend to recur.
\end{proof}

The theorem has an empirical side. If crises recur, there should be identifiable episodes beyond Husserl's, with the structure he described: a successful idealizing method, its technization, its transmission without reactivation, the substitution of its objects for the life-world, and a loss of the capacity to say what its results mean. We consider three.

\section{First episode: Galileo's mathematization of nature}\label{s:galileo}

Husserl's own episode is the template. Galileo's claim that the book of nature is written in the language of mathematics \citep{Galileo1623} expresses the program in its most compressed form. The program presupposed geometry as a sedimented technique: Galileo did not reactivate the sense of geometrical idealization, which he inherited from a tradition two thousand years old, and he applied it to new material, the motions and qualities of bodies. The application had three steps, each of which Husserl analyzes. The shapes of bodies were replaced by geometrical limit-shapes; the plenum of sensible qualities was treated as an index of quantitative processes, so that color became a matter of motion; and the resulting manifold was described by formulae, which allowed predictions to be computed without reference to what the formulae described.

The step from idealization to substitution occurs in the distinction between primary and secondary qualities. Galileo held that colors, tastes and odors exist only in the sensing body, while figure, number and motion belong to bodies themselves \citep{Galileo1623}. On Husserl's reading, this distinction expresses the forgetting directly: what the method could not mathematize directly was declared subjective, and the world from which the method drew its sense was redescribed as the effect, in us, of the mathematical world. \citet{Koyre1957} described the result as the destruction of the cosmos and the geometrization of space, and Koyré's history of the Galilean revolution, which Husserl's text anticipates in part, shows how the new physics replaced the qualitative world of Aristotelian science with the homogeneous space of geometry.

We accept Husserl's analysis, with one qualification. Husserl writes as though the substitution were complete in Galileo. It was not: Galileo still argued by appeal to everyday experience, and the full positivist self-interpretation that Husserl attacks took three centuries to develop. The Galilean episode is better described as the beginning of a crisis than as one already present, and this supports the recurrence thesis. If the crisis took time to develop, it was a process rather than an event, and processes of the same kind can begin again.

\section{Second episode: the institutionalization of statistical inference}\label{s:statistics}

The second episode is the rise of statistical inference in the social and life sciences between the early nineteenth and the late twentieth century. Its history has been written by \citet{Hacking1990} and \citet{Porter1995}, and it displays the structure of Theorem~\ref{t:recurrence} with unusual clarity.

The founding operations are familiar. Statistics begins with the counting of births, deaths and marriages in the registers of states and parishes: a life-world practice of record-keeping. Quetelet's \emph{homme moyen}, the average man, is an idealization of these counts: a being whose characteristics are the means of the distributions of characteristics in a population \citep{Quetelet1835}. Quetelet treated the average man as a real type, of which actual individuals are deviations, and the step from idealization to substitution is visible in his claim that the average man is the true object of social physics, while individuals are errors around him. Hacking describes the subsequent history as the \enquote{taming of chance}: the transformation of statistical regularities into laws governing populations, and the extension of the normal curve from errors of measurement to the characteristics of persons \citep{Hacking1990}.

The second stage is technization. In the twentieth century, Fisher's methods of significance testing \citep{Fisher1925}, together with the Neyman--Pearson theory of hypothesis testing, were combined in textbooks into a hybrid procedure that could be applied without understanding either of its sources. \citet{Gigerenzer2004} has called the result a \enquote{null ritual}: the researcher sets up a null hypothesis of no effect, tests it at a conventional level, and reports a result as significant if the null is rejected. The ritual was taught to generations of researchers as a technique, and most of its users could not have reconstructed the reasoning that gives it sense. Gigerenzer documents surveys in which large majorities of researchers and even statistics instructors misinterpreted the meaning of a $p$-value. This is Principle~\ref{p:transmission} in its purest form: a method used correctly, in the sense of being applied according to its rules, by practitioners who cannot reactivate its sense.

The third stage is substitution. In fields that adopted the ritual, the statistically significant effect came to be treated as the object of research, and the question that occasioned the research, whether a phenomenon in the life-world was real and how large it was, receded. An effect that was significant was taken to exist, and its existence was then reported as a finding about human beings, even when the effect was too small to be detected in anyone's experience, or when the significance was produced by the flexibility of the analysis rather than by the phenomenon.

The fourth stage is crisis. The replication crisis in psychology and related fields, which became visible in the 2010s when large-scale replication projects failed to reproduce a substantial share of published findings \citep{OSC2015}, is the crisis in the sense of Definition~\ref{d:crisis}. It was not, in the first instance, a crisis of error: the individual errors were familiar and had been criticized for decades. It was a crisis of self-understanding, in which fields discovered that they could not say what their accumulated results meant, because the results had been produced by a technique whose relation to the phenomena it was meant to detect had been forgotten. The response has included the reactivation of the sense of statistical inference: the reform movement's insistence on effect sizes, preregistration and the distinction between exploratory and confirmatory research is, in Husserl's terms, a return from the sedimented technique to the questions that gave it sense.

The episode coheres with the account of science defended elsewhere in the collection. SCI-01 identifies science by institutionalized error-correction, and SCI-19 holds that science is self-correcting, but slowly, through replication and meta-analysis that correct the record with a substantial delay. The replication crisis is an instance of slow self-correction, and the present account says why the correction was slow. Error-correction operates within a founded extension: it detects errors in results by the standards of the method. A crisis is not an error in results but a forgetting of what the method is for, and it is detected only when the accumulated results are confronted with the life-world from which their questions arose. That confrontation, in the replication crisis, took the form of a simple question, whether the effects could be seen again, which the method had made it possible not to ask.

\section{Third episode: computational science}\label{s:computation}

The third episode is in progress. Since the mid-twentieth century, a growing share of scientific knowledge has been produced by computation: simulations of systems too complex for analytic solution, and, more recently, machine learning models that are fitted to data and used for prediction without an explicit theory of the system they predict. \citet{Humphreys2004} argued that computational science extends human epistemic capacities beyond their natural range, and \citet{Humphreys2009} that simulation is epistemically novel because it is \emph{epistemically opaque}: no human can survey all the steps by which a simulation reaches its results.

Epistemic opacity is technization carried to its limit. In the Galilean case, the formulae could be applied without reactivating their sense, but the practitioner could in principle follow each step of the calculation. In the case of a large simulation, the practitioner cannot follow the steps, and in the case of a trained model, there may be no steps whose sense could be reactivated, since the parameters of the model were fitted to data and correspond to no idealization that anyone performed. \citet{Burrell2016} distinguishes three sources of opacity in machine learning: intentional secrecy, technical illiteracy and the opacity that arises from the scale and structure of the models themselves. Only the third is new, and it is the one that matters here: it makes reactivation impossible, not merely uncommon.

The founding operations are nonetheless present. A simulation is constructed from a model, which is an idealization in Husserl's sense: SCI-10 holds that models are false but represent through partial similarity, and that their idealizations are essential to their explanatory power. A machine learning model is trained on data, which are records of observations or of human judgements made in the life-world, and evaluated against further such records. The sense of its outputs derives from the sense of the data, and its evidence from the comparison of its outputs with life-world observations. Computational science is thus a founded extension in the sense of Definition~\ref{d:founded}, and Theorem~\ref{t:priority} applies to it.

The signs of forgetting are already visible. When a model's predictions are treated as findings about the world without the question being asked what in the data they reflect, the substitution of Definition~\ref{d:forgetting} occurs. Cases in which classifiers learned to detect features of the image-acquisition process rather than the medical conditions they were meant to diagnose, or in which models trained on historical decisions reproduced the biases of those decisions while being reported as neutral, have the structure of the replication crisis in compressed form: a technique applied correctly by its own standards yields results whose relation to the phenomena has been lost. Whether the episode will develop into a crisis in the full sense of Definition~\ref{d:crisis} depends on whether the field institutionalizes reactivation, that is, the practice of tracing a model's outputs back to the data and the life-world questions that give them sense. The recurrence thesis does not predict the outcome, but it predicts the tendency.

\section{Further episodes}\label{s:further}

Two further episodes can be noted briefly, as they illustrate the generality of Theorem~\ref{t:recurrence} beyond the natural and social sciences.

The first is the crisis in the foundations of mathematics in the early twentieth century, which Weyl named the \enquote{new foundational crisis} \citep{Weyl1921}. The paradoxes of set theory arose in a mathematics that had extended its techniques, the formation of arbitrary sets and the use of impredicative definitions, beyond the domain in which their sense had been constituted. Weyl's response, and Brouwer's, can be read as attempts at reactivation: a return to the constructions from which mathematical sense derives. Husserl himself, in \emph{Formal and Transcendental Logic}, described the formal mathematics of his time as a technique whose sense required clarification. We note the episode without adopting the intuitionist response, which MTH-07 rejects: reactivation does not require restricting mathematics to what can be constructed, only recovering the sense of what it does.

The second is economics. SCI-14 holds that economics is an idealized science whose models apply to actual agents only approximately. The history of the discipline includes repeated episodes in which the idealized agent of the models was taken for the real one, and in which the failure of the models to anticipate events in the economy that people live in prompted a crisis of self-understanding. The structure is that of Theorem~\ref{t:recurrence}: idealization, technization in the form of standard models, transmission without reactivation, substitution and crisis.

Each episode differs from Husserl's in its material, and each has the same form. This is what the recurrence thesis claims.

\section{Founded extension and scientific realism}\label{s:realism}

A reader may suspect that the priority thesis is anti-realist: that if science is founded on the life-world, its objects are constructions, and the life-world is the only reality. Husserl's own formulation, which speaks of the life-world as the \enquote{only real world}, encourages the suspicion. We argue that the account defended here is compatible with scientific realism, and in particular with SCI-02, on which mature, predictively successful theories are approximately true.

The compatibility rests on the distinction drawn after Theorem~\ref{t:priority}, between the priority of sense and evidence and a priority of being. Theorem~\ref{t:priority} says that the concepts of science have their sense through idealization of life-world practices and that their evidence is finally given in the life-world. It does not say that what the concepts apply to is constituted by the idealization. The concept of an electron derives its sense from a history of experimental practices and theoretical idealizations, and the evidence for electron theory is given in the readings of instruments; electrons exist independently of both. The forgetting of Definition~\ref{d:forgetting} is not the belief that electrons exist. It is the belief that the world consists \emph{only} of what physics describes in its idealized terms, and that the colored, textured, meaningful world of experience is an illusion produced by it. Realism about electrons does not require that belief, and the present account rejects it.

Read in this way, Husserl's surreptitious substitution is a mistake about the relation between theory and world, not a mistake in theory. The realist says that mature theories are approximately true of the world. The phenomenologist adds that the world of which they are true is the world given in the life-world, that the theories describe it under idealizations whose sense derives from life-world practices, and that the theories are not on that account descriptions of a different, more real world behind the one we experience. The two claims are consistent. SCI-10, on which models are false but represent through partial similarity, makes the same point at the level of models: a model represents the world by idealizing it, and to take the idealization for the world is to mistake the model for its target.

Two further theses of the collection's philosophy of science fit the account. SCI-17 holds that the effectiveness of mathematics in science is not unreasonable, because scientists select the mathematics that fits the phenomena. This is the Galilean episode described from the other side: idealizations are chosen for their fit with the life-world phenomena they are to predict, and their effectiveness is the effectiveness of a well-chosen instrument. SCI-20 holds that a final theory probably exists, though it may exceed human comprehension. A final theory would be the completion of the founded extension that physics is, and its existence would not abolish the priority of the life-world, since its concepts would still have their sense and its evidence through the operations and observations from which it was reached. If it exceeds human comprehension, it is the limiting case of Principle~\ref{p:transmission}: a theory whose sense no human could reactivate. Such a theory could be true. It would also be maximally exposed to the forgetting described here, since nothing in the practice of using it would hold its relation to the life-world in view. The recurrence thesis thus identifies a risk that grows with the success of science, without casting doubt on the truth of its results.

\section{Objections and replies}

\begin{objection}[Husserl's crisis was unique]
Husserl presents the crisis as a crisis of European humanity, in which the idea of reason that has guided European history since the Greeks is threatened. A crisis of that scope happens once. The episodes of Sections~\ref{s:statistics} to~\ref{s:further} are crises of particular disciplines, not of reason.
\end{objection}

\begin{reply}
We grant that the teleological framing of the \emph{Crisis} makes the crisis unique, and we do not adopt that framing. As \citet{Carr1974} argued, the analysis of sedimentation that the \emph{Crisis} contains does not depend on the teleology, and Husserl's own account of the Galilean episode identifies a mechanism, the technization and transmission of a method, that is not specific to the seventeenth century. Our thesis concerns the mechanism. If the crisis of European humanity was the cumulative effect of many disciplinary crises together with a general self-interpretation of culture by science, then it was itself an instance of the recurrent structure, at the largest scale.
\end{reply}

\begin{objection}[Forgetting is not a defect]
\citet{Blumenberg1981} argued that technization is the condition of efficiency: a technique whose foundations had to be reactivated at every use would be no technique. If forgetting is necessary, it is not a crisis.
\end{objection}

\begin{reply}
Blumenberg is right that forgetting in the sense of non-reactivation is necessary, and Proposition~\ref{p:tendency} accordingly states a tendency that cannot be eliminated. The crisis is not non-reactivation as such but the loss of the capacity for it, together with the substitution of Definition~\ref{d:forgetting}. A science in which most practitioners do not reactivate but some can, and in which the relation of the method to the life-world can be recovered when needed, is not in crisis. The replication crisis shows the difference: the techniques had been used without reactivation for decades without crisis, and the crisis arose when it emerged that no one could say what the accumulated results meant.
\end{reply}

\begin{objection}[The life-world is itself theory-laden]
The contemporary life-world is saturated with science: we see the sunset as the rotation of the earth, and we understand illness in terms of germs. There is no pre-scientific life-world on which science could be founded.
\end{objection}

\begin{reply}
Husserl anticipated the objection: the results of science flow back into the life-world and become part of it \citep[§34]{Husserl1954}. But the life-world on which a science is founded is not the life-world as it was before all science; it is the world of the practices from which the science's concepts are idealized and in which its observations are made. That world may contain sedimented results of earlier science; what matters is that the experiences in which they are given, such as seeing a thermometer's reading, are not themselves the idealized objects of the science founded on them. The theory-ladenness of the life-world thus confirms the account: it is the sedimentation of earlier founded extensions, which are available to perception only in the form in which they have become self-evident. \citet{Moran2012} notes the same layering in Husserl's own descriptions of the life-world.
\end{reply}

\begin{objection}[Realism and the only real world]
Husserl calls the life-world the only real world. If the account is compatible with realism, it is not Husserl's account.
\end{objection}

\begin{reply}
We do not claim that Husserl was a scientific realist in the contemporary sense; his transcendental idealism makes the question complex. We claim that the priority thesis and the forgetting thesis, as stated in Definitions~\ref{d:founded} and~\ref{d:forgetting}, are compatible with realism, and that Husserl's polemical formula is best read as a rejection of the substitution, not as a denial that scientific entities exist. \citet{Gurwitsch1974} and \citet{Sellars1962}, from different directions, show that a priority of the life-world in the order of sense is compatible with a priority of scientific entities in the order of being.
\end{reply}

\begin{objection}[Reactivation is impossible]
Derrida argued, in his introduction to the \emph{Origin of Geometry}, that full reactivation is impossible, since the origin of geometry is accessible only through the traces it has left, and that the ideal of reactivation is therefore an illusion \citep{Derrida1962}. If so, every science is always already in crisis.
\end{objection}

\begin{reply}
Reactivation in Definition~\ref{d:sedimentation} does not require recovering the historical acts of the first geometers. It requires re-performing operations of the same kind and so making the sense of the concepts evident again. The reform of statistical practice is an example: researchers who learn to ask what a $p$-value means, by reconstructing the reasoning of a test, reactivate its sense without recovering Fisher's thoughts. Derrida's point applies to historical origins; the account needs only constitutive ones.
\end{reply}

\section{What the result does not establish}

The result does not establish that any particular science is now in crisis, nor that crises are inevitable: Proposition~\ref{p:tendency} states a tendency that institutions of reflection can counteract. It does not establish that the life-world is ontologically prior to the entities of science, a claim that Section~\ref{s:realism} rejects. It does not establish that the content of scientific theories should be revised in the light of phenomenological reflection; the remedy for crisis is reactivation of sense, not revision of results. And it does not settle the larger question of Husserl's philosophy of history, whether European reason has a teleology, which the recurrence thesis neither requires nor excludes. Finally, the episodes are offered as evidence for the recurrence thesis, not as complete historical analyses; each has been the subject of a large literature from which we have drawn only what bears on the structure of Theorem~\ref{t:recurrence}.

\section{Conclusion}

The life-world is prior to science in the order of sense and evidence: every science idealizes life-world practices and confirms its results in life-world experience, and so is a founded extension of the life-world. Founded extensions tend to forget their foundation, because their methods are technized and transmitted to practitioners who use them without reactivating their sense, and who come to take the idealized objects of the method for the whole of reality. The crisis that Husserl diagnosed in Galilean physics is therefore recurrent, and its structure can be found in the institutionalization of statistical inference, in the opacity of computational science and in other episodes. The priority of the life-world is compatible with scientific realism, since it concerns the sense of scientific concepts and not the existence of what they describe; the crisis is a misunderstanding of what science is, and its remedy is reflection, not revision.

\begin{thebibliography}{99}\small
\bibitem[Blumenberg(1981)]{Blumenberg1981} Blumenberg, H. (1981). Lebenswelt und Technisierung unter Aspekten der Phänomenologie. In \emph{Wirklichkeiten in denen wir leben}. Reclam. First published 1963.
\bibitem[Blumenberg(2010)]{Blumenberg2010} Blumenberg, H. (2010). \emph{Theorie der Lebenswelt}, ed. M.~Sommer. Suhrkamp.
\bibitem[Burrell(2016)]{Burrell2016} Burrell, J. (2016). How the machine \enquote{thinks}: Understanding opacity in machine learning algorithms. \emph{Big Data \& Society}, 3(1).
\bibitem[Carr(1974)]{Carr1974} Carr, D. (1974). \emph{Phenomenology and the Problem of History: A Study of Husserl's Transcendental Philosophy}. Northwestern University Press.
\bibitem[Derrida(1962)]{Derrida1962} Derrida, J. (1962). Introduction. In E.~Husserl, \emph{L'origine de la géométrie}, trans. J.~Derrida. Presses Universitaires de France.
\bibitem[Fisher(1925)]{Fisher1925} Fisher, R.~A. (1925). \emph{Statistical Methods for Research Workers}. Oliver and Boyd.
\bibitem[Galilei(1623)]{Galileo1623} Galilei, G. (1623). \emph{Il Saggiatore}. Mascardi.
\bibitem[Gigerenzer(2004)]{Gigerenzer2004} Gigerenzer, G. (2004). Mindless statistics. \emph{The Journal of Socio-Economics}, 33(5), 587--606.
\bibitem[Gurwitsch(1974)]{Gurwitsch1974} Gurwitsch, A. (1974). \emph{Phenomenology and the Theory of Science}, ed. L.~Embree. Northwestern University Press.
\bibitem[Hacking(1990)]{Hacking1990} Hacking, I. (1990). \emph{The Taming of Chance}. Cambridge University Press.
\bibitem[Heelan(1983)]{Heelan1983} Heelan, P.~A. (1983). \emph{Space-Perception and the Philosophy of Science}. University of California Press.
\bibitem[Humphreys(2004)]{Humphreys2004} Humphreys, P. (2004). \emph{Extending Ourselves: Computational Science, Empiricism, and Scientific Method}. Oxford University Press.
\bibitem[Humphreys(2009)]{Humphreys2009} Humphreys, P. (2009). The philosophical novelty of computer simulation methods. \emph{Synthese}, 169(3), 615--626.
\bibitem[Husserl(1939)]{Husserl1939b} Husserl, E. (1939). Die Frage nach dem Ursprung der Geometrie als intentional-historisches Problem, ed. E.~Fink. \emph{Revue internationale de philosophie}, 1(2), 203--225.
\bibitem[Husserl(1954)]{Husserl1954} Husserl, E. (1954). \emph{Die Krisis der europäischen Wissenschaften und die transzendentale Phänomenologie}, ed. W.~Biemel. Husserliana VI. Martinus Nijhoff. English translation by D.~Carr, \emph{The Crisis of European Sciences and Transcendental Phenomenology}, Northwestern University Press, 1970.
\bibitem[Koyré(1957)]{Koyre1957} Koyré, A. (1957). \emph{From the Closed World to the Infinite Universe}. Johns Hopkins Press.
\bibitem[Moran(2012)]{Moran2012} Moran, D. (2012). \emph{Husserl's Crisis of the European Sciences and Transcendental Phenomenology: An Introduction}. Cambridge University Press.
\bibitem[Open Science Collaboration(2015)]{OSC2015} Open Science Collaboration (2015). Estimating the reproducibility of psychological science. \emph{Science}, 349(6251), aac4716.
\bibitem[Porter(1995)]{Porter1995} Porter, T.~M. (1995). \emph{Trust in Numbers: The Pursuit of Objectivity in Science and Public Life}. Princeton University Press.
\bibitem[Quetelet(1835)]{Quetelet1835} Quetelet, A. (1835). \emph{Sur l'homme et le développement de ses facultés, ou Essai de physique sociale}. Bachelier.
\bibitem[Sellars(1962)]{Sellars1962} Sellars, W. (1962). Philosophy and the scientific image of man. In R.~Colodny (Ed.), \emph{Frontiers of Science and Philosophy}, 35--78. University of Pittsburgh Press.
\bibitem[Weyl(1921)]{Weyl1921} Weyl, H. (1921). Über die neue Grundlagenkrise der Mathematik. \emph{Mathematische Zeitschrift}, 10, 39--79.
\end{thebibliography}

\end{document}
